\documentclass[10pt,letterpaper]{article}
\usepackage{cornell-notes}

\title{Clause 33.08.03: Class Template Counting Semaphore}
\author{}
\date{}
\setDocTitle{Clause 33.08.03: Class Template Counting Semaphore}
\setDocSubtitle{C++ 2024}
\setDocOwner{C++ 2024 Cornell Notes}
\setDocDate{\today}
\setCornellCollection{C++ 2024 Cornell Notes}
\setCornellUnitType{Clause}
\setCornellUnitNumber{33.08.03}
\setCornellUnitTitle{Class Template Counting Semaphore}
\hypersetup{pdftitle={Clause 33.08.03: Class Template Counting Semaphore},pdfsubject={Cornell notes},pdfauthor={}}

\begin{document}
\maketitle
\begin{CornellOverviewBox}{Overview}
\textbf{Classification:} Standard library - Normative\\[2pt]
\textbf{Learning objective:} Understand the rules, terminology, and semantic consequences associated with Class template counting\_semaphore.
\end{CornellOverviewBox}\begin{CornellSummaryBox}{Summary}
{\sffamily\bfseries\color{Accent} Summary}\par\vspace{3pt}Section 33.8.3 focuses on Class template counting\_semaphore. Apply the specified interface together with its mandates, effects, result conditions, exception behavior, and complexity constraints. When applying it, preserve the distinction between mandatory requirements, recommendations, permissions, and behavior deliberately left undefined, unspecified, or implementation-defined.
\end{CornellSummaryBox}

\end{document}
